\begin{textsample}{POS Dim 1 – 2009 – Score 29.00 – 2009/t000143.txt}  \label{ex:f1_pos_029}
Fifty years ago, when I began exploring \textbf{the} ocean, no one — not Jacques Perrin, not Jacques Cousteau or Rachel Carson — imagined that we could do anything to harm \textbf{the} ocean by what we put into it or by what we took out of it.
It seemed, at that time, to be a \textbf{sea} of Eden, but now we know, and now we are facing paradise lost.
I want to share with you my personal view of changes in \textbf{the} \textbf{sea} that affect all of us, and to consider why it matters that in 50 years, we ' ve lost — actually, we ' ve taken, we ' ve eaten — more than 90 percent of \textbf{the} big \textbf{fish} in \textbf{the} \textbf{sea} ; why you should care that nearly half of \textbf{the} coral reefs have \textbf{disappeared} ; why a mysterious depletion of oxygen in large areas of \textbf{the} Pacific should concern not only \textbf{the} creatures that are dying, but it really should concern you.
It does concern you, as well.
I ' m haunted by \textbf{the} thought of what Ray Anderson calls"tomorrow ' s child,"asking why we didn ' t do something on our watch to save sharks and bluefin tuna and squids and coral reefs and \textbf{the} living ocean while there still was time.
Well, now is that time.
I hope for your help to explore and protect \textbf{the} wild ocean in ways that will restore \textbf{the} health and, in so doing, secure hope for humankind.
Health to \textbf{the} ocean means health for us.
And I hope Jill Tarter ' s wish to engage Earthlings includes dolphins and \textbf{whales} and other \textbf{sea} creatures in this quest to find intelligent life elsewhere in \textbf{the} \textbf{universe}.
And I hope, Jill, that someday we will find evidence that there is intelligent life among humans on this \textbf{planet}.
Did I say that?
I guess I did.
For me, as a scientist, it all began in 1953 when I first tried scuba.
It ' s when I first got to know \textbf{fish} swimming in something other than lemon slices and butter.
I actually love diving at night ; you see a lot of \textbf{fish} then that you don ' t see in \textbf{the} daytime.
Diving day and night was really easy for me in 1970, when I led a team of aquanauts living \textbf{underwater} for weeks at a time — at \textbf{the} same time that astronauts were putting their footprints on \textbf{the} moon.
In 1979 I had a chance to put my footprints on \textbf{the} ocean floor while using this personal submersible called Jim.
It was six miles offshore and 1, 250 feet down.
It ' s one of my favorite bathing suits.
Since then, I ' ve used about 30 kinds of submarines and I ' ve started three companies and a nonprofit foundation called Deep Search to design and build systems to access \textbf{the} deep \textbf{sea}.
I led a five-year National Geographic expedition, \textbf{the} Sustainable \textbf{Seas} expeditions, using these little subs.
They ' re so simple to drive that even a scientist can do it.
And I ' m living proof.
Astronauts and aquanauts alike really appreciate \textbf{the} importance of air, food, water, temperature — all \textbf{the} things you need to stay alive in space or under \textbf{the} \textbf{sea}.
I heard astronaut Joe Allen explain how he had to learn everything he could about his life support system and then do everything he could to take care of his life support system ; and then he pointed to this and he said,"Life support system."We need to learn everything we can about it and do everything we can to take care of it. \textbf{The} poet Auden said,"Thousands have lived without love ; none without water."Ninety-seven percent of Earth ' s water is ocean.
No \textbf{blue}, no green.
If you think \textbf{the} ocean isn ' t important, imagine Earth without it.
Mars comes to mind.
No ocean, no life support system.
I gave a talk not so long ago at \textbf{the} World Bank and I showed this amazing image of Earth and I said,"There it is! \textbf{The} World Bank!"That ' s where all \textbf{the} assets are!
And we ' ve been trawling them down much faster than \textbf{the} natural systems can replenish them.
Tim Worth says \textbf{the} economy is a wholly-owned subsidiary of \textbf{the} environment.
With every drop of water you \textbf{drink}, every breath you take, you ' re connected to \textbf{the} \textbf{sea}.
No matter where on Earth you live.
Most of \textbf{the} oxygen in \textbf{the} atmosphere is generated by \textbf{the} \textbf{sea}.
Over time, most of \textbf{the} \textbf{planet} ' s \textbf{organic} \textbf{carbon} has been absorbed and stored there, mostly by microbes. \textbf{The} ocean drives climate and weather, stabilizes temperature, shapes Earth ' s \textbf{chemistry}.
Water from \textbf{the} \textbf{sea} forms clouds that return to \textbf{the} land and \textbf{the} \textbf{seas} as rain, sleet and snow, and provides home for about 97 percent of life in \textbf{the} world, maybe in \textbf{the} \textbf{universe}.
No water, no life ; no \textbf{blue}, no green.
Yet we have this idea, we humans, that \textbf{the} Earth — all of it : \textbf{the} oceans, \textbf{the} skies — are so vast and so resilient it doesn ' t matter what we do to it.
That may have been true 10, 000 years ago, and maybe even 1, 000 years ago but in \textbf{the} last 100, especially in \textbf{the} last 50, we ' ve drawn down \textbf{the} assets, \textbf{the} air, \textbf{the} water, \textbf{the} wildlife that make our lives possible.
New technologies are helping us to understand \textbf{the} nature of nature ; \textbf{the} nature of what ' s happening, showing us our impact on \textbf{the} Earth.
I mean, first you have to know that you ' ve got a problem.
And fortunately, in our time, we ' ve learned more about \textbf{the} problems than in all preceding history.
And with knowing comes caring.
And with caring, there ' s hope that we can find an enduring place for ourselves within \textbf{the} natural systems that support us.
But first we have to know.
Three years ago, I met John Hanke, who ' s \textbf{the} head of Google Earth, and I told him how much I loved being able to hold \textbf{the} world in my hands and go exploring vicariously.
But I asked him :"When are you going to finish it?
You did a great job with \textbf{the} land, \textbf{the} dirt.
What about \textbf{the} water?"Since then, I ' ve had \textbf{the} great pleasure of working with \textbf{the} Googlers, with DOER Marine, with National Geographic, with dozens of \textbf{the} best institutions and scientists around \textbf{the} world, ones that we could enlist, to put \textbf{the} ocean in Google Earth.
And as of just this week, last Monday, Google Earth is now whole.
Consider this : Starting right here at \textbf{the} convention center, we can find \textbf{the} nearby aquarium, we can look at where we ' re sitting, and then we can cruise up \textbf{the} coast to \textbf{the} big aquarium, \textbf{the} ocean, and California ' s four national marine sanctuaries, and \textbf{the} new network of state marine reserves that are beginning to protect and restore some of \textbf{the} assets We can flit over to Hawaii and see \textbf{the} real Hawaiian Islands : not just \textbf{the} little bit that pokes through \textbf{the} \textbf{surface}, but also what ' s below.
To see — wait a minute, we can go kshhplash! — right there, ha — under \textbf{the} ocean, see what \textbf{the} \textbf{whales} see.
We can go explore \textbf{the} other side of \textbf{the} Hawaiian Islands.
We can go actually and swim around on Google Earth and visit with humpback \textbf{whales}.
These are \textbf{the} gentle giants that I ' ve had \textbf{the} pleasure of meeting face to face many times \textbf{underwater}.
There ' s nothing quite like being personally inspected by a \textbf{whale}.
We can pick up and fly to \textbf{the} deepest place : seven miles down, \textbf{the} Mariana Trench, where only two people have ever been.
Imagine that.
It ' s only seven miles, but only two people have been there, 49 years ago.
One-way trips are easy.
We need new deep-diving submarines.
How about some X Prizes for ocean exploration?
We need to see deep trenches, \textbf{the} undersea mountains, and understand life in \textbf{the} deep \textbf{sea}.
We can now go to \textbf{the} Arctic.
Just ten years ago I stood on \textbf{the} ice at \textbf{the} North Pole.
An ice-free Arctic Ocean may happen in this century.
That ' s bad news for \textbf{the} polar bears.
That ' s bad news for us too.
Excess \textbf{carbon} dioxide is not only driving global warming, it ' s also changing ocean \textbf{chemistry}, making \textbf{the} \textbf{sea} more acidic.
That ' s bad news for coral reefs and oxygen-producing plankton.
Also it ' s bad news for us.
We ' re putting hundreds of millions of tons of plastic and other trash into \textbf{the} \textbf{sea}.
Millions of tons of \textbf{discarded} \textbf{fishing} nets, gear that continues to kill.
We ' re clogging \textbf{the} ocean, poisoning \textbf{the} \textbf{planet} ' s circulatory system, and we ' re taking out hundreds of millions of tons of wildlife, all carbon-based units.
Barbarically, we ' re killing sharks for shark \textbf{fin} \textbf{soup}, undermining food chains that shape planetary \textbf{chemistry} and drive \textbf{the} \textbf{carbon} cycle, \textbf{the} nitrogen cycle, \textbf{the} oxygen cycle, \textbf{the} water cycle — our life support system.
We ' re still killing bluefin tuna ; truly endangered and much more valuable alive than dead.
All of these parts are part of our life support system.
We kill using long lines, with baited hooks every few feet that may stretch for 50 miles or more.
Industrial trawlers and draggers are scraping \textbf{the} \textbf{sea} floor like bulldozers, taking everything in their path.
Using Google Earth you can witness trawlers — in China, \textbf{the} North Sea, \textbf{the} Gulf of Mexico — shaking \textbf{the} foundation of our life support system, leaving plumes of death in their path. \textbf{The} next time you dine on sushi — or sashimi, or swordfish steak, or \textbf{shrimp} cocktail, whatever wildlife you happen to \textbf{enjoy} from \textbf{the} ocean — think of \textbf{the} real cost.
For every \textbf{pound} that goes to market, more than 10 \textbf{pounds}, even 100 \textbf{pounds}, may be thrown away as bycatch.
This is \textbf{the} consequence of not knowing that there are limits to what we can take out of \textbf{the} \textbf{sea}.
This chart shows \textbf{the} decline in ocean wildlife from 1900 to 2000. \textbf{The} highest concentrations are in red.
In my lifetime, imagine, 90 percent of \textbf{the} big \textbf{fish} have been killed.
Most of \textbf{the} \textbf{turtles}, sharks, tunas and \textbf{whales} are way down in numbers.
But, there is good news.
Ten percent of \textbf{the} big \textbf{fish} still remain.
There are still some \textbf{blue} \textbf{whales}.
There are still some krill in Antarctica.
There are a few oysters in Chesapeake Bay.
Half \textbf{the} coral reefs are still in pretty good shape, a jeweled belt around \textbf{the} middle of \textbf{the} \textbf{planet}.
There ' s still time, but not a lot, to turn things around.
But business as usual means that in 50 years, there may be no coral reefs — and no commercial \textbf{fishing}, because \textbf{the} \textbf{fish} will simply be gone.
Imagine \textbf{the} ocean without \textbf{fish}.
Imagine what that means to our life support system.
Natural systems on \textbf{the} land are in big trouble too, but \textbf{the} problems are more obvious, and some actions are being taken to protect trees, watersheds and wildlife.
And in 1872, with Yellowstone National Park, \textbf{the} United States began establishing a system of parks that some say was \textbf{the} best idea America ever had.
About 12 percent of \textbf{the} land around \textbf{the} world is now protected : safeguarding biodiversity, providing a \textbf{carbon} sink, generating oxygen, protecting watersheds.
And, in 1972, this nation began to establish a counterpart in \textbf{the} \textbf{sea}, National Marine Sanctuaries.
That ' s another great idea. \textbf{The} good news is that there are now more than 4, 000 places in \textbf{the} \textbf{sea}, around \textbf{the} world, that have some kind of protection.
And you can find them on Google Earth. \textbf{The} bad news is that you have to look hard to find them.
In \textbf{the} last three years, for example, \textbf{the} U.
S. protected 340, 000 square miles of ocean as national monuments.
But it only increased from 0. 6 of one percent to 0. 8 of one percent of \textbf{the} ocean protected, globally.
Protected areas do rebound, but it takes a long time to restore 50-year-old rockfish or monkfish, sharks or \textbf{sea} bass, or 200-year-old orange roughy.
We don ' t consume 200-year-old cows or chickens.
Protected areas provide hope that \textbf{the} creatures of Ed Wilson ' s dream of an encyclopedia of life, or \textbf{the} census of marine life, will live not just as a list, a photograph, or a paragraph.
With scientists around \textbf{the} world, I ' ve been looking at \textbf{the} 99 percent of \textbf{the} ocean that is open to \textbf{fishing} — and mining, and drilling, and dumping, and whatever — to \textbf{search} out hope spots, and try to find ways to give them and us a secure future.
Such as \textbf{the} Arctic — we have one chance, right now, to get it right.
Or \textbf{the} Antarctic, where \textbf{the} \textbf{continent} is protected, but \textbf{the} surrounding ocean is being stripped of its krill, \textbf{whales} and \textbf{fish}.
Sargasso Sea ' s three million square miles of floating forest is being gathered up to feed cows. 97 percent of \textbf{the} land in \textbf{the} Galapagos Islands is protected, but \textbf{the} adjacent \textbf{sea} is being ravaged by \textbf{fishing}.
It ' s true too in Argentina on \textbf{the} Patagonian shelf, which is now in serious trouble. \textbf{The} high \textbf{seas}, where \textbf{whales}, tuna and dolphins travel — \textbf{the} largest, least protected, ecosystem on Earth, filled with luminous creatures, living in dark waters that average two miles deep.
They flash, and sparkle, and glow with their own living light.
There are still places in \textbf{the} \textbf{sea} as pristine as I knew as a child. \textbf{The} next 10 years may be \textbf{the} most important, and \textbf{the} next 10, 000 years \textbf{the} best chance our species will have to protect what remains of \textbf{the} natural systems that give us life.
To cope with climate change, we need new ways to generate power.
We need new ways, better ways, to cope with poverty, wars and disease.
We need many things to keep and maintain \textbf{the} world as a better place.
But, nothing else will matter if we fail to protect \textbf{the} ocean.
Our fate and \textbf{the} ocean ' s are one.
We need to do for \textbf{the} ocean what Al Gore did for \textbf{the} skies above.
A global plan of action with a world conservation union, \textbf{the} IUCN, is underway to protect biodiversity, to mitigate and recover from \textbf{the} impacts of climate change, on \textbf{the} high \textbf{seas} and in coastal areas, wherever we can identify critical places.
New technologies are needed to map, photograph and explore \textbf{the} 95 percent of \textbf{the} ocean that we have yet to see. \textbf{The} goal is to protect biodiversity, to provide stability and resilience.
We need deep-diving subs, new technologies to explore \textbf{the} ocean.
We need, maybe, an expedition — a TED at \textbf{sea} — that could help figure out \textbf{the} next steps.
And so, I suppose you want to know what my wish is.
I wish you would use all means at your disposal — films, expeditions, \textbf{the} \textbf{web}, new submarines — and campaign to ignite public support for a global network of marine protected areas — hope spots large enough to save and restore \textbf{the} ocean, \textbf{the} \textbf{blue} heart of \textbf{the} \textbf{planet}.
How much?
Some say 10 percent, some say 30 percent.
You decide : how much of your heart do you want to protect?
Whatever it is, a \textbf{fraction} of one percent is not enough.
My wish is a big wish, but if we can make it happen, it can truly change \textbf{the} world, and help ensure \textbf{the} survival of what actually — as it turns out — is my favorite species ; that would be us.
For \textbf{the} children of today, for tomorrow ' s child : as never again, now is \textbf{the} time.
Thank you.

% matched lemmas: blue, carbon, chemistry, continent, disappear, discard, drink, enjoy, fin, fish, fishing, fraction, organic, planet, pound, sea, search, shrimp, soup, surface, the, turtle, underwater, universe, web, whale
\end{textsample}
